
% ----------------------------------------------------------------------
%                   LATEX TEMPLATE FOR PhD THESIS
% ----------------------------------------------------------------------

% based on Harish Bhanderi's PhD/MPhil template, then Uni Cambridge
% http://www-h.eng.cam.ac.uk/help/tpl/textprocessing/ThesisStyle/
% corrected and extended in 2007 by Jakob Suckale, then MPI-CBG PhD programme
% and made available through OpenWetWare.org - the free biology wiki


%: Style file for Latex
% Most style definitions are in the external file PhDthesisPSnPDF.
% In this template package, it can be found in ./Latex/Classes/
\documentclass[a4paper, twoside,11pt, openright]{Latex/Classes/PhDthesisPSnPDF}

\usepackage[polish]{babel}
\usepackage{polski}
\usepackage{url}	
\usepackage{amssymb,amsmath}
\usepackage{verbatim} 
\usepackage{amsmath}
\usepackage{array}
\usepackage{graphics}
\usepackage{subfigure}
\usepackage[utf8]{inputenc}
%% for code listings
\usepackage{listings}
\lstset{language=Java,captionpos=b,tabsize=3,frame=lines,keywordstyle=\color{black},commentstyle=\color{black},stringstyle=\color{black},breaklines=true,showstringspaces=false,basicstyle=\footnotesize,aboveskip={1.5\baselineskip},belowskip={1.5\baselineskip}}


%: Macro file for Latex
% Macros help you summarise frequently repeated Latex commands.
% Here, they are placed in an external file /Latex/Macros/MacroFile1.tex
% An macro that you may use frequently is the figuremacro (see introduction.tex)
% This file contains macros that can be called up from connected TeX files
% It helps to summarise repeated code, e.g. figure insertion (see below).

% insert a centered figure with caption and description
% parameters 1:filename, 2:title, 3:description and label
\newcommand{\figuremacro}[3]{
	\begin{figure}[htbp]
		\centering
		\includegraphics[width=1\textwidth]{#1}
		\caption[#2]{\textbf{#2} - #3}
		\label{#1}
	\end{figure}
}

% insert a centered figure with caption and description AND WIDTH
% parameters 1:filename, 2:title, 3:description and label, 4: textwidth
% textwidth 1 means as text, 0.5 means half the width of the text
\newcommand{\figuremacroW}[4]{
	\begin{figure}[htbp]
		\centering
		\includegraphics[width=#4\textwidth]{#1}
		\caption[#2]{\textbf{#2} - #3}
		\label{#1}
	\end{figure}
}

% inserts a figure with wrapped around text; only suitable for NARROW figs
% o is for outside on a double paged document; others: l, r, i(inside)
% text and figure will each be half of the document width
% note: long captions often crash with adjacent content; take care
% in general: above 2 macro produce more reliable layout
\newcommand{\figuremacroN}[3]{
	\begin{wrapfigure}{o}{0.5\textwidth}
		\centering
		\includegraphics[width=0.48\textwidth]{#1}
		\caption[#2]{{\small\textbf{#2} - #3}}
		\label{#1}
	\end{wrapfigure}
}

% predefined commands by Harish
\newcommand{\PdfPsText}[2]{
  \ifpdf
     #1
  \else
     #2
  \fi
}

\newcommand{\IncludeGraphicsH}[3]{
  \PdfPsText{\includegraphics[height=#2]{#1}}{\includegraphics[bb = #3, height=#2]{#1}}
}

\newcommand{\IncludeGraphicsW}[3]{
  \PdfPsText{\includegraphics[width=#2]{#1}}{\includegraphics[bb = #3, width=#2]{#1}}
}

\newcommand{\InsertFig}[3]{
  \begin{figure}[!htbp]
    \begin{center}
      \leavevmode
      #1
      \caption{#2}
      \label{#3}
    \end{center}
  \end{figure}
}


%%% Local Variables: 
%%% mode: latex
%%% TeX-master: "~/Documents/LaTeX/CUEDThesisPSnPDF/thesis"
%%% End: 


%: --------------------------------------------------------------
%:                  FRONT MATTER: dedications, abstract,..
% --------------------------------------------------------------

\begin{document}


\def\tablename{Tabela}
%\def\prefacename{Przedmowa}%
%\def\refname{Literatura}%
%\def\abstractname{Streszczenie}%
%\def\bibname{Bibliografia}%
%\def\chaptername{Rozdzia\l}%
%\def\appendixname{Dodatek}%
%\def\contentsname{Spistre\'sci}% spis treści
%\def\listfigurename{Spisrysunk\'ow}% spis rysunków
\def\listtablename{Spis tabel}% spis tablic
%\def\indexname{Indeks}%
%\def\figurename{Rysunek}%
%\def\tablename{Tablica}%
%\def\partname{Cz\eob{}\'s\'c}% część
%\def\enclname{Za\l\aob{}cznik}% załącznik
%\def\ccname{Kopie:}%
%\def\headtoname{Do}%
%\def\pagename{Strona}%
%\def\seename{Por\'ownaj}% porównaj
%\def\alsoname{Por\'ownajtak\.ze}% porównaj także
%\def\proofname{Dow\'od}% dowód
%\def\glossaryname{Glossary}

%\language{polish}

% sets line spacing
\renewcommand\baselinestretch{1.2}
\baselineskip=18pt plus1pt


%: ----------------------- generate cover page ------------------------


\thispagestyle{empty}

\begin{center}
Akademia Górniczo - Hutnicza im. Stanisława Staszica w Krakowie \\
Wydział Elektrotechniki, Automatyki, Informatyki i Elektroniki \\
Katedra Informatyki
% \\ Akademii Górniczo Hutniczej w Krakowie
\rule{\textwidth}{.1mm}
\end{center}

\begin{figure}[hkp!]
 \centering
% \includegraphics[width=50pt]{agh_znak}
 \includegraphics[bb=0 0 64 124]{logo}
 \label{fig:orzel_agh}
\end{figure}

\begin{center}
{\bf \large Tytuł pracy}
\vskip 10pt
{\bf \Large Integracja systemów przetwarzania mowy w środowisku zgodnym z paradygmatem SOA}
\end{center}

\vskip 60pt
\begin{center}
{\Large \bf Konrad Dziedzic, Rafał Fronczyk}\\
{rfronczyk@gmail.com, konraddziedzic@gmail.com}
\end{center}

\vskip 50pt
\begin{center}
{\bf
    Promotor pracy: dr inż. Łukasz Czekierda
}
\end{center}

\noindent
\rule{\textwidth}{.1mm}
\begin{center}
Kraków, 2012
\end{center}


%: ----------------------- oswiadczenie -------------------


\clearpage
\thispagestyle{empty}

\begin{center}
  \textbf{Oświadczenie}
\end{center}

\vskip 20pt

Oświadczamy, świadomi odpowiedzialności karnej za poświadczenie nieprawdy, że niniejszą pracę dyplomową wykonaliśmy osobiście i samodzielnie (w zakresie wyszczególnionym we wstępie) i że nie korzystaliśmy ze źródeł innych niż wymienione w pracy.

\vskip 60pt

\begin{center}
\begin{tabular}{c c}

..................................................... & ..................................................... \\
\begin{footnotesize}
\textit{podpis i data}
\end{footnotesize}
&
\begin{footnotesize}
\textit{podpis i data}
\end{footnotesize}
\\

\end{tabular}
\end{center}

\clearpage



%: ----------------------- abstract ------------------------

% Your institution may have specific regulations if you need an abstract and where it is to be placed in the document. The default here is just after title.


% Thesis Abstract -----------------------------------------------------

%\begin{abstractslong}    %uncommenting this line, gives a different abstract heading
\begin{abstracts}        %this creates the heading for the abstract page

Rozwój technologiczny, a w szczególności rozwój komputerów i technologii z nimi związanych w ciągu ostatnich lat jest bardzo szybki. Zmieniają się zarówno, podzespoły, moc obliczeniowa, rozmiary, wygląd a nawet interfejsy. Komputery wkroczyły, lub są temu bardzo bliskie, do niemal każdej dziedziny życia. W związku z tym zmienia się sposób komunikacji między użytkownikiem a komputerem. W ostatnich czasach można zaobserwować dążenie inżynierów i projektantów do uczynienia komunikacji z komputerem jak najbardziej naturalną. Pomysły są różne od "kontrolerów bez kontrolera" jak Microsoft Kinect \footnote {http://www.xbox.com/en-US/kinect}, poprzez rozwiązania "ruchowe" podobne do tego na jakie zdecydowało się Sony \footnote{http://www.sony.com/} w swoim kontrolerze PlayStation Move \footnote{http://us.playstation.com/ps3/playstation-move/} poprzez klasyczne  jak mysz i klawiatura. Jednak najbardziej naturalnym sposobem porozumiewania się wydaje się głos. Najbardziej znanym, bo napewno nie pionierskim, systemem który komunikuję się z użytkownikiem za pomocą głosu jest Apple Siri \footnote{http://www.apple.com/iphone/features/siri.html} - osobisty asystent, "żyjacy" wewnątrz systemu, umożliwiający dostęp do jego funkcji za pomocą mowy.\\
Niniejsza praca przedstawia prototyp systemu UniversalSynthesizer, służacego do zamiany tekstu na dźwięk i mowy na tekst, będącego zaawansowaną, rozproszoną, rozszerzalną, wielojęzyczną platformą/serwisem umożliwiającą łatwe tworzenie zróżnicowanych, wieloplatformowych aplikacji, mających różne zadania. Celem pracy nie było stworzenie gotowego do użytku, kompletnego, w pełni sprawnego produktu. Powstałą aplikację należy traktować bardziej jako punkt wyjścia, prototyp który w przyszłości może być wykorzystany do zbudowanie w pełni funkcjonalnego, ogólnie dostępnego, komercyjnego lub otwartego systemu. W związku z tym niektóre zagadnienia, dość istotne z punktu widzenia potencjalnego odbiorcy, ale nie będące ściśle powiązane z celem pracy zostały pominięty lub też niedopracowane. 
\end{abstracts}
%\end{abstractlongs}


% ---------------------------------------------------------------------- 


% The original template provides and abstractseparate environment, if your institution requires them to be separate. I think it's easier to print the abstract from the complete thesis by restricting printing to the relevant page.
% \begin{abstractseparate}
%   
% Thesis Abstract -----------------------------------------------------

%\begin{abstractslong}    %uncommenting this line, gives a different abstract heading
\begin{abstracts}        %this creates the heading for the abstract page

Rozwój technologiczny, a w szczególności rozwój komputerów i technologii z nimi związanych w ciągu ostatnich lat jest bardzo szybki. Zmieniają się zarówno, podzespoły, moc obliczeniowa, rozmiary, wygląd a nawet interfejsy. Komputery wkroczyły, lub są temu bardzo bliskie, do niemal każdej dziedziny życia. W związku z tym zmienia się sposób komunikacji między użytkownikiem a komputerem. W ostatnich czasach można zaobserwować dążenie inżynierów i projektantów do uczynienia komunikacji z komputerem jak najbardziej naturalną. Pomysły są różne od "kontrolerów bez kontrolera" jak Microsoft Kinect \footnote {http://www.xbox.com/en-US/kinect}, poprzez rozwiązania "ruchowe" podobne do tego na jakie zdecydowało się Sony \footnote{http://www.sony.com/} w swoim kontrolerze PlayStation Move \footnote{http://us.playstation.com/ps3/playstation-move/} poprzez klasyczne  jak mysz i klawiatura. Jednak najbardziej naturalnym sposobem porozumiewania się wydaje się głos. Najbardziej znanym, bo napewno nie pionierskim, systemem który komunikuję się z użytkownikiem za pomocą głosu jest Apple Siri \footnote{http://www.apple.com/iphone/features/siri.html} - osobisty asystent, "żyjacy" wewnątrz systemu, umożliwiający dostęp do jego funkcji za pomocą mowy.\\
Niniejsza praca przedstawia prototyp systemu UniversalSynthesizer, służacego do zamiany tekstu na dźwięk i mowy na tekst, będącego zaawansowaną, rozproszoną, rozszerzalną, wielojęzyczną platformą/serwisem umożliwiającą łatwe tworzenie zróżnicowanych, wieloplatformowych aplikacji, mających różne zadania. Celem pracy nie było stworzenie gotowego do użytku, kompletnego, w pełni sprawnego produktu. Powstałą aplikację należy traktować bardziej jako punkt wyjścia, prototyp który w przyszłości może być wykorzystany do zbudowanie w pełni funkcjonalnego, ogólnie dostępnego, komercyjnego lub otwartego systemu. W związku z tym niektóre zagadnienia, dość istotne z punktu widzenia potencjalnego odbiorcy, ale nie będące ściśle powiązane z celem pracy zostały pominięty lub też niedopracowane. 
\end{abstracts}
%\end{abstractlongs}


% ---------------------------------------------------------------------- 

% \end{abstractseparate}


%: ----------------------- contents ------------------------

\setcounter{secnumdepth}{3} % organisational level that receives a numbers
\setcounter{tocdepth}{3}    % print table of contents for level 3
\tableofcontents            % print the table of contents
% levels are: 0 - chapter, 1 - section, 2 - subsection, 3 - subsection



%: --------------------------------------------------------------
%:                  MAIN DOCUMENT SECTION
% --------------------------------------------------------------

% the main text starts here with the introduction, 1st chapter,...
\mainmatter

\renewcommand{\chaptername}{} % uncomment to print only "1" not "Chapter 1"


%: ----------------------- subdocuments ------------------------

% Parts of the thesis are included below. Rename the files as required.
% But take care that the paths match. You can also change the order of appearance by moving the include commands.


% this file is called up by thesis.tex
% content in this file will be fed into the main document

%: ----------------------- introduction file header -----------------------
\chapter{Wstęp}

% the code below specifies where the figures are stored
\ifpdf
    \graphicspath{{1_introduction/figures/PNG/}{1_introduction/figures/PDF/}{1_introduction/figures/}}
\else
    \graphicspath{{1_introduction/figures/EPS/}{1_introduction/figures/}}
\fi

% ----------------------------------------------------------------------
%: ----------------------- introduction content ----------------------- 
% ----------------------------------------------------------------------



%: ----------------------- HELP: latex document organisation
% the commands below help you to subdivide and organise your thesis
%    \chapter{}       = level 1, top level
%    \section{}       = level 2
%    \subsection{}    = level 3
%    \subsubsection{} = level 4
% note that everything after the percentage sign is hidden from output



\section{Definicja problemu } % section headings are printed smaller than chapter names
% intro
Integracja to łączenie mniejszych elementów ze sobą w celu stworzenia większego, spójnego, poprawnie działającego elementu. W informatyce integracja dotyczy głównie całych systemów. Czasami integracja nie kończy się tylko na łączeniu funkcjonalności ale także na dodawaniu nowych, jest to spowodowane tym, że system mający wiele komponentów mających różne zastosowania ma dzięki temu dużo większe możliwości i możę uzyskać funkcjonalność jakiej jego składowe nie byłyby w stanie uzyskać samodzielnie. Integracja systemów jest zadaniem bardzo trudnym, mającym charakter przekrojowy, wymagającyn wiedzy z wielu dziedzin takich jak:
 \begin{itemize}
	\item sieci komputerowe
	\item programowanie
	\item zarządzanie procesami biznesowymi
	\item znajomość narzędzi wspierających integrację
\end{itemize} 
Komputerowe przetwarzanie mowy jest w ostatnim czasie szybko rozwijającą się dziedziną w przemyśle informatycznym. Wykorzystuje się je w wielu różnych dziedzinach zarówno związanych z rozrywką ( rozmainte gry polegające na jak najwierniejszym zaśpiewaniu podanego utworu), z życiem codziennym (Apple Siri czyli wirtualny asystent) jak i z pracą(różne systemy telefoniczne działające bez udziału człowieka). W wielu sytuacjach wykorzystanie komputerowo przetwarzanj mowy jest niezastąpione.\\
Podobnie sytuacja wygląda z wykorzystaniem architektury SOA. Jest to jedno z najważniejszych osiągnieć informatycznych ostatnich lat. Podejście to zdecydowanie dominuje w obecnie projektowanych architekturach. Jednym z głowych powodów takiego stanu rzeczy jest uproszczenie modelu współczesnych dużych aplikacji. Umożliwia ono podzielenie systemu na serwisy które są niezależnymi, samodzielnymi, łatwymi w zarządzaniu, posiadającymi jasno i jednoznacznie zdefiniowane interfejsy encjami, które można wykorzystać w różnej kolejności w celu otrzymania różnych efektów. \\
Celem tej pracy jest zbadanie możliwości wykorzystania nowoczesnej i prężnie rozwijającej się architektury do stworzenia gotowego systemy ułatwiającego i ujednolicającego proces tworzenia aplikacji przetwarzających mowę.
 


\subsection{Obszar badań} % subsection headings are again smaller than section names
W skład zakresu pracy wchodzą zarówno technologie związane z architekturą SOA jak i z przetwarzaniem mowy. W celu użycia odpowiednich narzędzi przeprowadzono wiele badań i testów w czasie których dużą wagę przykładano do:
\begin{itemize}
 	\item możliwości konfiguracyjnych
	\item wydajności
	\item łatwości rozbudowy
	\item dostępnej dokumentacjji i pomocy technicznej
\end{itemize}
W efekcie przeprowadzonych czynności okazało się, że najlepsze efekty uzyskać można korzystając z szyny ESB  Apache ServiceMix,  syntezatorów mowy Ivona i FreeTTS oraz innych mniej istotnych narzędzi potrzebnych do stworzenia architektury zgodnej z SOA oraz zdolnej do efektywnego wykonywania postawionych przed nią zadań. 


	% background information
% this file is called up by thesis.tex
% content in this file will be fed into the main document

\chapter{Aims of the project} % top level followed by section, subsection


% ----------------------- contents from here ------------------------

\section{Final aim}

Our ultimate goal is...

\section{Preliminary aims}

There will be several preliminary scientific targets to be accomplished on the way...




% ---------------------------------------------------------------------------
% ----------------------- end of thesis sub-document ------------------------
% ---------------------------------------------------------------------------						% aims of the project

\include{3/XYZ}			
\include{4/XYZ}	
\include{5/XYZ}
\include{6/XYZ}

% this file is called up by thesis.tex
% content in this file will be fed into the main document

\chapter{Podsumowanie} % top level followed by section, subsection


% ----------------------- paths to graphics ------------------------

% change according to folder and file names
\ifpdf
    \graphicspath{{7/figures/PNG/}{7/figures/PDF/}{7/figures/}}
\else
    \graphicspath{{7/figures/EPS/}{7/figures/}}
\fi

Obecnie istnieje bardzo dużo róznych serwisów oferujących przetwarzanie mowy. Jednak ich poważną wadą jest to, że dziedzina ich pracy jest bardzo wąska. Brakuje systemów oferujących więcej niż jeden sposób transformacji. Autorzy tej pracy starali się sprawdzic czy integracja takich systemów jest dobrym pomysłem oraz jeżeli tak to jakie podejście jest najlepsze. Jak zostało to zaprezentowane w powyższych rozdziałach stworzenie systemu integrującego serwisy posiadające różne funkcje ma dużo zalet i bardzo ułatwia tworzenie aplikacji wykorzystujących te serwisy. Co więcej dodatkowa warstwa abstrakcji, która powstaje w ten sposób, nie wpływa w sposób znaczący na szybkość działania aplikacji klienckich (w porównaniu z sytuacją w której aplikacja kliencka sama musiałaby komunikować się z tymi serwisami). Wszystko to wyraźnie pokazuje, że istnienie takiego zintegrowanego serwisu ma sens i jest na niego zapotrzebowanie. \\
Co do wyboru sposobu integracji, jak zostało to zaprezentowane wyżej istnieje wiele odrębnych sposobów, z których każdy ma swoje zalety i wady. Po dogłębnej analizie dostępnych rozwiązań integracyjnych autorzy pracy zdecydowali, że najlepszy efekt daje najnowsze rozwiązanie a więc ESB. Posiada ono szereg zalet które należy uwypuklić:
\begin{itemize}
	\item Łatwość rozbudowy - gdy powstanie jakiś nowy serwis, oferujący ciekawą funkcjonalność jedynym wysiłkiem jaki należy wykonać w celu użycia go, jest wpięcie serwisu do szyny
	\item Lekkość - warstwa abtrakcji nakładana przez ESB w celu połączenia dostępnych serwisów jest bardzo mała i nie ma wpływu na wydajność
	\item Skalowalność - różne instancje ESB mogą łączyć się razem w celu uzyskania wyższej wydajności oraz rozproszenia geograficznego
	\item Różnorodność - wszystkie istniejące implementacje ESB oferują wiele różnych endpoint'ów, służących jako interfejsy wejścia/wyjścia, obsługujących wiele formatów
\end{itemize}
Wykorzystanie takiego podejścia daje też inny duży, plus jakim niewątpliwie jest łatwość ukrywania implementacji. Klienci systemu opartego na takim rozwiązaniu nie będą wiedzieć z jakich konkretnie serwisów korzystają, przez co w razie wystąpienia kłopotów czy też pojawienia się nowszego, lepszego lub ciekawszego rozwiązania można taki serwis łatwo podmienić, bez konieczności zmieniania API. Kolejną dużą zaletą wykorzystania ESB, jest możliwość wykorzystania routingu. Dzięki temu sterowanie całym systemem i kolejnością wywoływania poszczególnych serwisów jest możliwa za pomocą prostego pliku xml.  \\
Zaprezentowana praca magisterska zagłębia się w dziedzinę, która jest dość nowatorska. Zaproponowane rozwiązanie i przykładowa implementacja pokazują zarówno, iż temat jest ciekawy i warty dalszych badań oraz, że istniejące w tej chwili rozwiązania integracyjne są odpowiednie dla specyficznej dziedziny jaką jest przetwarzanie mowy i można je z powodzeniem stosować. 

% ----------------------- contents from here ------------------------






% ---------------------------------------------------------------------------
% ----------------------- end of thesis sub-document ------------------------
% ---------------------------------------------------------------------------               % discussion of results


% this file is called up by thesis.tex
% content in this file will be fed into the main document

\chapter{Materials \& methods} % top level followed by section, subsection


% ----------------------- paths to graphics ------------------------

% change according to folder and file names
\ifpdf
    \graphicspath{{8/figures/PNG/}{8/figures/PDF/}{8/figures/}}
\else
    \graphicspath{{8/figures/EPS/}{8/figures/}}
\fi

% ----------------------- contents from here ------------------------




 

% ---------------------------------------------------------------------------
%: ----------------------- end of thesis sub-document ------------------------
% ---------------------------------------------------------------------------



 






        % description of lab methods




% --------------------------------------------------------------
%:                  BACK MATTER: appendices, refs,..
% --------------------------------------------------------------

% the back matter: appendix and references close the thesis


%: ----------------------- bibliography ------------------------

% The section below defines how references are listed and formatted
% The default below is 2 columns, small font, complete author names.
% Entries are also linked back to the page number in the text and to external URL if provided in the BibTex file.

% PhDbiblio-url2 = names small caps, title bold & hyperlinked, link to page 
%\begin{multicols}{1} % \begin{multicols}{ # columns}[ header text][ space]
%\begin{tiny} % tiny(5) < scriptsize(7) < footnotesize(8) < small (9)

\bibliographystyle{Latex/Classes/PhDbiblio-url2} % Title is link if provided
\renewcommand{\bibname}{Bibliografia} % changes the header; default: Bibliography

\bibliography{9_backmatter/references} % adjust this to fit your BibTex file

%\end{tiny}
%\end{multicols}

% --------------------------------------------------------------
% Various bibliography styles exit. Replace above style as desired.

% in-text refs: (1) (1; 2)
% ref list: alphabetical; author(s) in small caps; initials last name; page(s)
%\bibliographystyle{Latex/Classes/PhDbiblio-case} % title forced lower case
%\bibliographystyle{Latex/Classes/PhDbiblio-bold} % title as in bibtex but bold
%\bibliographystyle{Latex/Classes/PhDbiblio-url} % bold + www link if provided

%\bibliographystyle{Latex/Classes/jmb} % calls style file jmb.bst
% in-text refs: author (year) without brackets
% ref list: alphabetical; author(s) in normal font; last name, initials; page(s)

%\bibliographystyle{plainnat} % calls style file plainnat.bst
% in-text refs: author (year) without brackets
% (this works with package natbib)


% --------------------------------------------------------------

% according to Dresden med fac summary has to be at the end
%
% Thesis Abstract -----------------------------------------------------

%\begin{abstractslong}    %uncommenting this line, gives a different abstract heading
\begin{abstracts}        %this creates the heading for the abstract page

Rozwój technologiczny, a w szczególności rozwój komputerów i technologii z nimi związanych w ciągu ostatnich lat jest bardzo szybki. Zmieniają się zarówno, podzespoły, moc obliczeniowa, rozmiary, wygląd a nawet interfejsy. Komputery wkroczyły, lub są temu bardzo bliskie, do niemal każdej dziedziny życia. W związku z tym zmienia się sposób komunikacji między użytkownikiem a komputerem. W ostatnich czasach można zaobserwować dążenie inżynierów i projektantów do uczynienia komunikacji z komputerem jak najbardziej naturalną. Pomysły są różne od "kontrolerów bez kontrolera" jak Microsoft Kinect \footnote {http://www.xbox.com/en-US/kinect}, poprzez rozwiązania "ruchowe" podobne do tego na jakie zdecydowało się Sony \footnote{http://www.sony.com/} w swoim kontrolerze PlayStation Move \footnote{http://us.playstation.com/ps3/playstation-move/} poprzez klasyczne  jak mysz i klawiatura. Jednak najbardziej naturalnym sposobem porozumiewania się wydaje się głos. Najbardziej znanym, bo napewno nie pionierskim, systemem który komunikuję się z użytkownikiem za pomocą głosu jest Apple Siri \footnote{http://www.apple.com/iphone/features/siri.html} - osobisty asystent, "żyjacy" wewnątrz systemu, umożliwiający dostęp do jego funkcji za pomocą mowy.\\
Niniejsza praca przedstawia prototyp systemu UniversalSynthesizer, służacego do zamiany tekstu na dźwięk i mowy na tekst, będącego zaawansowaną, rozproszoną, rozszerzalną, wielojęzyczną platformą/serwisem umożliwiającą łatwe tworzenie zróżnicowanych, wieloplatformowych aplikacji, mających różne zadania. Celem pracy nie było stworzenie gotowego do użytku, kompletnego, w pełni sprawnego produktu. Powstałą aplikację należy traktować bardziej jako punkt wyjścia, prototyp który w przyszłości może być wykorzystany do zbudowanie w pełni funkcjonalnego, ogólnie dostępnego, komercyjnego lub otwartego systemu. W związku z tym niektóre zagadnienia, dość istotne z punktu widzenia potencjalnego odbiorcy, ale nie będące ściśle powiązane z celem pracy zostały pominięty lub też niedopracowane. 
\end{abstracts}
%\end{abstractlongs}


% ---------------------------------------------------------------------- 



%: ----------------------- list of figures/tables ------------------------

\listoffigures	% print list of figures

\listoftables  % print list of tables

\end{document}
